\documentclass[11pt,a4paper]{article}
\usepackage[utf8]{inputenc}
\usepackage[T1]{fontenc}
\usepackage{geometry}
\usepackage{graphicx}
\usepackage{hyperref}
\usepackage{listings}
\usepackage{xcolor}
\usepackage{amsmath}
\usepackage{tikz}
\usepackage{float}
\usepackage{enumitem}

% Page setup
\geometry{margin=1in}
\hypersetup{
    colorlinks=true,
    linkcolor=blue,
    filecolor=magenta,      
    urlcolor=cyan,
    pdftitle={SPAM Code Structure},
    pdfauthor={SPAM Development Team}
}

% Code listing style
\lstset{
    language=Python,
    basicstyle=\ttfamily\small,
    keywordstyle=\color{blue}\bfseries,
    commentstyle=\color{green!60!black},
    stringstyle=\color{red},
    numbers=left,
    numberstyle=\tiny\color{gray},
    stepnumber=1,
    numbersep=5pt,
    backgroundcolor=\color{gray!10},
    frame=single,
    breaklines=true,
    breakatwhitespace=true,
    tabsize=4,
    showspaces=false,
    showstringspaces=false
}

\title{\textbf{SPAM Code Structure Documentation}\\
\large Scanner for Polarized Anisotropic Materials\\
\large Code Architecture Overview}
\author{SPAM Development Team}
\date{\today}

\begin{document}

\maketitle
\tableofcontents
\newpage

\section{Project Overview}

The SPAM (Scanner for Polarized Anisotropic Materials) project is a Python-based desktop application for automated material characterization. The codebase is organized into a modular structure with clear separation between GUI, database, and hardware control components.

\section{Project Structure}

\subsection{Directory Layout}
\begin{verbatim}
SPAM/
├── GUI.py                    # Main application file
├── motor_control_status.py   # Motor control reference
├── requirements.txt          # Python dependencies
├── start_spam.sh            # Linux startup script
├── start_spam.bat           # Windows startup script
├── spam_config.json         # Configuration file (auto-generated)
├── spam.db                  # SQLite database (auto-generated)
├── backend/
│   ├── __init__.py          # Package marker
│   ├── database.py          # Database configuration
│   └── models.py            # Database models
└── venv/                    # Virtual environment
\end{verbatim}

\subsection{Key Files}

\begin{description}
    \item[\texttt{GUI.py}] Main application file containing the entire GUI and application logic (3000+ lines)
    \item[\texttt{backend/database.py}] Database connection and session management
    \item[\texttt{backend/models.py}] SQLAlchemy ORM models for data storage
    \item[\texttt{motor\_control\_status.py}] Reference implementation for motor control
    \item[\texttt{requirements.txt}] Python package dependencies
    \item[\texttt{start\_spam.sh/bat}] Platform-specific startup scripts
\end{description}

\section{Main Components}

\subsection{GUI.py - Main Application}

\texttt{GUI.py} is the core of the application, containing a single main class that handles all functionality.

\subsubsection{SPAMGui Class}
The \texttt{SPAMGui} class inherits from \texttt{tk.Tk}, making it the main application window.

\textbf{Key Responsibilities:}
\begin{itemize}
    \item GUI component creation and layout
    \item User interaction handling
    \item Measurement workflow management
    \item Motor control integration
    \item Database operations
    \item Real-time data visualization
\end{itemize}

\textbf{Initialization Flow:}
\begin{enumerate}
    \item Window setup (size, title, colors)
    \item State variable initialization
    \item Load configuration from JSON file
    \item Create all UI components
    \item Initialize database connection
    \item Initialize motor control (if on Raspberry Pi)
    \item Start periodic display updates
\end{enumerate}

\subsubsection{Main Methods by Category}

\textbf{UI Creation Methods:}
\begin{itemize}
    \item \texttt{\_create\_menu()} - Creates menu bar
    \item \texttt{\_create\_sidebar()} - Creates left control panel
    \item \texttt{\_create\_info\_panel()} - Creates right info display
    \item \texttt{\_create\_center\_panel()} - Creates graph area
    \item \texttt{\_create\_control\_panel()} - Creates bottom control panel
    \item \texttt{\_create\_status\_bar()} - Creates status bar
\end{itemize}

\textbf{Measurement Methods:}
\begin{itemize}
    \item \texttt{\_on\_start\_measurement()} - Starts measurement thread
    \item \texttt{\_on\_stop\_measurement()} - Stops measurement
    \item \texttt{\_measurement\_worker()} - Background measurement thread
    \item \texttt{\_on\_calibrate()} - Two-step calibration process
\end{itemize}

\textbf{Motor Control Methods:}
\begin{itemize}
    \item \texttt{\_initialize\_motor\_control()} - Initializes I2C/GPIO
    \item \texttt{\_send\_motor\_command()} - Sends I2C motor commands
    \item \texttt{\_wait\_for\_motor\_position()} - Waits for position reached
\end{itemize}

\textbf{Database Methods:}
\begin{itemize}
    \item \texttt{\_create\_measurement()} - Stores measurement in database
    \item \texttt{\_create\_calibration()} - Stores calibration record
    \item \texttt{\_get\_measurements()} - Retrieves measurements
\end{itemize}

\textbf{Display Methods:}
\begin{itemize}
    \item \texttt{\_update\_graphs()} - Updates all graphs
    \item \texttt{\_update\_display()} - Periodic display refresh
    \item \texttt{\_update\_status()} - Updates status bar
\end{itemize}

\textbf{Settings Methods:}
\begin{itemize}
    \item \texttt{\_on\_connection\_setup()} - Connection settings dialog
    \item \texttt{\_on\_adjust\_parameters()} - Parameter adjustment dialog
    \item \texttt{\_load\_connection\_settings()} - Loads config from JSON
    \item \texttt{\_save\_connection\_settings()} - Saves config to JSON
\end{itemize}

\subsubsection{Debug Console Class}
The \texttt{DebugConsole} class is a separate window for diagnostics:
\begin{itemize}
    \item System information tab
    \item Database information tab
    \item Measurement log tab
    \item Console log tab
    \item Configuration tab
    \item Motor control tab
\end{itemize}

\subsection{Backend Package}

\subsubsection{database.py}
Handles database connection and session management.

\textbf{Key Components:}
\begin{itemize}
    \item \texttt{SQLALCHEMY\_DATABASE\_URL} - SQLite database path
    \item \texttt{engine} - SQLAlchemy engine instance
    \item \texttt{SessionLocal} - Session factory for database operations
    \item \texttt{Base} - Base class for ORM models
\end{itemize}

\textbf{Purpose:}
\begin{itemize}
    \item Creates SQLite database file (\texttt{spam.db})
    \item Manages database connections
    \item Provides session factory for database operations
\end{itemize}

\subsubsection{models.py}
Defines database schema using SQLAlchemy ORM.

\textbf{Measurement Model:}
\begin{itemize}
    \item Stores angle sweep measurement data
    \item Fields: id, angle, permittivity, permeability, transmitted/reflected power/phase, timestamp
    \item Indexed on angle and timestamp for fast queries
\end{itemize}

\textbf{Calibration Model:}
\begin{itemize}
    \item Stores calibration records
    \item Fields: id, timestamp, parameters (JSON), status
    \item Used for two-step calibration process
\end{itemize}

\section{Data Flow}

\subsection{Application Startup}
\begin{enumerate}
    \item \texttt{main()} function called
    \item Creates \texttt{SPAMGui} instance
    \item \texttt{\_\_init\_\_()} initializes all components
    \item Database tables created (if first run)
    \item Configuration loaded from \texttt{spam\_config.json}
    \item Motor control initialized (if on Linux)
    \item GUI displayed and ready for user interaction
\end{enumerate}

\subsection{Measurement Workflow}
\begin{enumerate}
    \item User clicks "Start Measurement"
    \item \texttt{\_on\_start\_measurement()} called
    \item Measurement thread spawned (\texttt{\_measurement\_worker})
    \item Thread loops through angles (0° to 90°):
    \begin{enumerate}
        \item Send motor command via I2C
        \item Wait for position reached
        \item Read ADC signals (I/Q channels)
        \item Calculate power and phase
        \item Calculate material properties (ε, μ)
        \item Store measurement in database
        \item Update GUI graphs and displays
        \item Increment angle
    \end{enumerate}
    \item Thread completes, updates status
\end{enumerate}

\subsection{Data Storage Flow}
\begin{enumerate}
    \item Measurement data calculated
    \item \texttt{\_create\_measurement()} called
    \item Creates \texttt{Measurement} object with all fields
    \item Adds to database session
    \item Commits transaction
    \item Data persisted to \texttt{spam.db}
\end{enumerate}

\subsection{Display Update Flow}
\begin{enumerate}
    \item \texttt{\_update\_display()} called periodically (every 1 second)
    \item Calls \texttt{\_update\_graphs()}
    \item Queries database for latest measurements
    \item Extracts data arrays (angles, permittivity, permeability, power, phase)
    \item Updates Matplotlib graphs
    \item Updates info panel values
    \item Updates status indicators
    \item Schedules next update
\end{enumerate}

\section{Key Technologies}

\subsection{Python Libraries}
\begin{description}
    \item[\texttt{tkinter}] GUI framework (built-in)
    \item[\texttt{matplotlib}] Graph plotting and visualization
    \item[\texttt{sqlalchemy}] ORM for database operations
    \item[\texttt{numpy}] Numerical computations
    \item[\texttt{smbus}] I2C communication (Raspberry Pi)
    \item[\texttt{RPi.GPIO}] GPIO control (Raspberry Pi)
\end{description}

\subsection{Communication Protocols}
\begin{description}
    \item[\texttt{I2C}] Motor control communication (Raspberry Pi)
    \item[\texttt{GPIO}] Interrupt handling for motor status
    \item[\texttt{SQLite}] Local database storage
    \item[\texttt{JSON}] Configuration file format
\end{description}

\section{Configuration Management}

\subsection{Configuration File}
Settings stored in \texttt{spam\_config.json}:
\begin{itemize}
    \item I2C bus number
    \item ADC I2C address
    \item IF-I and IF-Q channel numbers
    \item Sampling rate
    \item Microcontroller I2C address
    \item ISR GPIO pin number
\end{itemize}

\subsection{Configuration Flow}
\begin{enumerate}
    \item On startup: Load settings from JSON file
    \item If file doesn't exist: Use default values
    \item User changes settings via dialog
    \item Settings validated
    \item Saved to JSON file
    \item Motor control reinitialized with new settings
\end{enumerate}

\section{Motor Control Integration}

\subsection{I2C Communication}
\begin{itemize}
    \item Uses \texttt{smbus} library on Raspberry Pi
    \item Communicates with microcontroller at address 0x55
    \item Sends 6-byte commands: \texttt{[command, motor\_num, float\_bytes...]}
    \item Position encoded as 4-byte little-endian float
\end{itemize}

\subsection{GPIO Interrupts}
\begin{itemize}
    \item ISR pin (default: GPIO 17) configured for interrupts
    \item Interrupt handler reads status register
    \item Detects collision (bit 0x01) and position reached (bit 0x02)
    \item Updates GUI via thread-safe \texttt{self.after()} calls
\end{itemize}

\subsection{Platform Detection}
\begin{itemize}
    \item Windows: Motor control in simulation mode
    \item Linux: Full hardware integration
    \item Automatic detection via \texttt{platform.system()}
\end{itemize}

\section{Threading Model}

\subsection{Main Thread}
\begin{itemize}
    \item GUI event loop (Tkinter)
    \item Handles user interactions
    \item Updates display
    \item Manages UI components
\end{itemize}

\subsection{Measurement Thread}
\begin{itemize}
    \item Background worker for measurements
    \item Runs \texttt{\_measurement\_worker()} function
    \item Performs motor control and data acquisition
    \item Updates GUI via \texttt{self.after()} for thread safety
\end{itemize}

\subsection{Thread Safety}
\begin{itemize}
    \item GUI updates from background thread use \texttt{self.after()}
    \item Database operations are thread-safe (SQLAlchemy)
    \item Motor control uses proper synchronization
\end{itemize}

\section{Graph Visualization}

\subsection{Graph Types}
The application displays four real-time graphs:
\begin{enumerate}
    \item \textbf{Permittivity vs Angle} - Material permittivity over sweep
    \item \textbf{Permeability vs Angle} - Material permeability over sweep
    \item \textbf{Power vs Angle} - Transmitted and reflected power
    \item \textbf{Phase vs Angle} - Transmitted and reflected phase
\end{enumerate}

\subsection{Graph Implementation}
\begin{itemize}
    \item Uses Matplotlib's \texttt{FigureCanvasTkAgg} for embedding
    \item Graphs contained in scrollable frame
    \item Auto-scaling based on data range
    \item Real-time updates from database queries
    \item Professional styling with color scheme
\end{itemize}

\section{Database Schema}

\subsection{Measurements Table}
\begin{verbatim}
CREATE TABLE measurements (
    id INTEGER PRIMARY KEY,
    angle FLOAT NOT NULL,
    permittivity FLOAT NOT NULL,
    permeability FLOAT NOT NULL,
    transmitted_power FLOAT,
    reflected_power FLOAT,
    transmitted_phase FLOAT,
    reflected_phase FLOAT,
    timestamp DATETIME NOT NULL
);
\end{verbatim}

\subsection{Calibrations Table}
\begin{verbatim}
CREATE TABLE calibrations (
    id INTEGER PRIMARY KEY,
    timestamp DATETIME NOT NULL,
    parameters JSON,
    status VARCHAR NOT NULL
);
\end{verbatim}

\section{Error Handling}

\subsection{Error Types}
\begin{itemize}
    \item Hardware errors (I2C, GPIO failures)
    \item Database errors (connection, query failures)
    \item Measurement errors (motor timeouts, collisions)
    \item User input errors (invalid values)
\end{itemize}

\subsection{Error Handling Strategy}
\begin{itemize}
    \item Try-except blocks around critical operations
    \item User-friendly error messages
    \item Logging to debug console
    \item Graceful degradation (continue when possible)
    \item Status bar updates for user feedback
\end{itemize}

\section{Export Functionality}

\subsection{Export Formats}
\begin{itemize}
    \item \textbf{CSV}: Comma-separated values for spreadsheet import
    \item \textbf{JSON}: Structured data format for programmatic access
\end{itemize}

\subsection{Export Process}
\begin{enumerate}
    \item User clicks "Export" button
    \item Selects format (CSV or JSON)
    \item Chooses save location
    \item All measurements retrieved from database
    \item Data formatted and written to file
    \item File saved with timestamp in filename
\end{enumerate}

\section{Startup Scripts}

\subsection{start\_spam.sh (Linux)}
\begin{itemize}
    \item Checks Python installation
    \item Creates virtual environment if needed
    \item Installs dependencies
    \item Handles DNS/network issues
    \item Sets display for VNC
    \item Launches GUI application
\end{itemize}

\subsection{start\_spam.bat (Windows)}
\begin{itemize}
    \item Checks Python installation
    \item Creates virtual environment if needed
    \item Installs dependencies
    \item Launches GUI application
\end{itemize}

\section{Code Organization Principles}

\subsection{Modularity}
\begin{itemize}
    \item Single main class (\texttt{SPAMGui}) for GUI
    \item Separate backend package for database
    \item Clear method naming conventions
    \item Logical grouping of related methods
\end{itemize}

\subsection{Separation of Concerns}
\begin{itemize}
    \item GUI logic in \texttt{GUI.py}
    \item Database logic in \texttt{backend/}
    \item Configuration separate from code
    \item Hardware abstraction for platform differences
\end{itemize}

\subsection{Maintainability}
\begin{itemize}
    \item Clear method names (\texttt{\_create\_}, \texttt{\_on\_}, \texttt{\_update\_})
    \item Comprehensive comments
    \item Consistent code style
    \item Error handling throughout
\end{itemize}

\section{Key Design Patterns}

\subsection{Singleton Pattern}
\begin{itemize}
    \item Single \texttt{SPAMGui} instance
    \item Single database session
    \item Single configuration file
\end{itemize}

\subsection{Observer Pattern}
\begin{itemize}
    \item GUI updates observe database changes
    \item Status updates observe measurement state
    \item Graph updates observe new measurements
\end{itemize}

\subsection{Factory Pattern}
\begin{itemize}
    \item \texttt{SessionLocal} creates database sessions
    \item Measurement objects created via factory method
\end{itemize}

\section{Summary}

The SPAM codebase follows a clear, modular structure:
\begin{itemize}
    \item \textbf{Main File}: \texttt{GUI.py} contains all GUI and application logic
    \item \textbf{Backend}: Database models and configuration in \texttt{backend/}
    \item \textbf{Configuration}: Settings stored in JSON file
    \item \textbf{Database}: SQLite for local data storage
    \item \textbf{Hardware}: I2C/GPIO integration for Raspberry Pi
    \item \textbf{Threading}: Background measurement thread with thread-safe GUI updates
    \item \textbf{Visualization}: Real-time graphs using Matplotlib
\end{itemize}

The code is designed for maintainability, with clear separation between GUI, database, and hardware control components. The modular structure makes it easy to understand, modify, and extend functionality.

\end{document}
